\chapter{Аналитический раздел}

В данном разделе представляются формализация задачи, описания загружаемых модулей Linux, известных характеристик использования процессора и оперативной памяти и анализ известных способов сбора данных об использовании системой процессора и оперативной памяти.

\section{Загружаемые модули ядра Linux}

Загружаемый модуль ядра представляет собой компилированный объектный код, предназначенный для динамического расширения функциональности ядра операционной системы без необходимости его перекомпиляции. Данный механизм позволяет инкапсулировать новый код в независимые компоненты, которые могут быть присоединены к исполняемому образу ядра во время его работы или отсоединены от него.

Процесс интеграции модуля заключается в вызове его функции инициализации, которая регистрирует предоставляемые модулем услуги в соответствующих подсистемах ядра. После успешной регистрации функция инициализации завершает работу, а модуль становится частью исполняемого кода ядра. Корректное удаление модуля инициирует выполнение функции завершения, задача которой -- осуществить дерегистрацию услуг и освободить все занятые системные ресурсы для обеспечения целостности ядра.

Важным аспектом архитектуры модулей является их зависимость от интерфейсов, экспортируемых самим ядром: модуль не может использовать функции, не объявленные в его публичном API. Это ограничение обеспечивает стабильность системы, но создает зависимость модуля от конкретной версии ядра. Основное преимущество использования загружаемых модулей заключается в значительном сокращении цикла разработки и отладки низкоуровневого кода, так как процесс итеративного тестирования не требует перезапуска всей системы.

\newpage

\section{Известные характеристики использования процессора}

Операционная система использует центральный процессор в качестве ключевого вычислительного ресурса для координации работы всех компонентов системы. Для оценки использования процессора может быть применен набор количественных и качественных характеристик.

\subsection{Процессорное время}

Процессорное время характеризует интервал времени, в течение которого процессор выполнял конкретную задачу. Операционная система классифицирует процессорное время, учет которого ведется для каждого процессорного ядра в отдельности. Классы процессорного времени представлены в таблице~\ref{tbl:processor-time-types}.

\begin{table}[H]
	\centering
	\caption{Классы процессорного времени}
	\label{tbl:processor-time-types}
	\begin{tabularx}{\textwidth}{|p{4cm}|X|}
		\hline
		\textbf{Класс} & \textbf{Описание} \\
		\hline
		Пользовательское время & Время, затраченное на исполнение кода программ в пространстве пользователя \\
		\hline
		Системное время & Время, затраченное на исполнение кода операционной системы в режиме ядра \\
		\hline
		Время ожидания ввода-вывода & Время, в течение которого процессор простаивал в состоянии готовности, ожидая завершения операций ввода-вывода от устройств \\
		\hline
		Время обработки прерываний & Время, затраченное на обслуживание аппаратных прерываний \\
		\hline
		Время обработки отложенных прерываний & Время, затраченное на обработку нижних половин прерываний \\
		\hline
		Время ожидания & Время, в течение которого процессорное ядро было полностью свободно \\
		\hline
	\end{tabularx}
\end{table}

\newpage

Для оценки относительной загрузки процессора $U_{CPU}$ за выбранный интервал используется показатель, вычисляемый как отношение времени полезной работы к полному времени:
\begin{equation}
	U_{CPU} = (1 - \frac{I_{end} - I_{begin}}{T_{end} - T_{begin}}) \times 100\%
\end{equation}
где:
\begin{enumerate}
	\item $I_{end}$ и $I_{begin}$ -- время ожидания в конце и начале интервала соответственно;
	\item $T_{end}$ и $T_{begin}$ -- полное время в конце и начале интервала соответственно.
\end{enumerate}

\subsection{Тактовая частота}

Такт ядра процессора -- промежуток между двумя импульсами тактового генератора, который синхронизирует выполнение всех операций процессора. Выполнение различных элементарных операций может занимать от долей такта до многих тактов в зависимости от команды и процессора. Тактовая частота является физической характеристикой процессора, определяющей количество синхронизирующих тактов, выполняемых за одну секунду. 

Современным процессорам свойственна способность к динамическому изменению тактовой частоты во время работы. Данный механизм служит для решения двух основных задач: повышения энергоэффективности и управления тепловыделением. Снижение частоты в периоды низкой вычислительной нагрузки позволяет значительно уменьшить энергопотребление и нагрев кристалла, а кратковременное повышение частоты сверх номинала позволяет увеличить производительность, если это позволяет тепловой пакет и система питания.

Ответственность за динамическое изменение тактовой частоты распределена между аппаратным обеспечением и программным обеспечением. Непосредственно генерация тактовой частоты осуществляется аппаратным блоком -- тактовым генератором, а логика принятия решения о необходимости и величине изменения частоты может реализовываться как аппаратным обеспечением, так и операционной системой.

\subsection{Количество инструкций за такт}

Количество инструкций за такт (IPC) характеризует эффективность использования процессором его вычислительных ресурсов в течение одного такта и может быть вычислено как:
\begin{equation}
	IPC = \frac{N_{instructions}}{N_{cycles}}
\end{equation}
где:
\begin{enumerate}
	\item $N_{instructions}$ -- общее количество исполненных процессором инструкций;
	\item $N_{cycles}$ -- общее количество затраченных процессорных тактов.
\end{enumerate}

Величина IPC не является постоянной, а изменяется динамически в зависимости от множества факторов. Наиболее значимое влияние оказывает микроархитектура процессора: количество и тип исполнительных устройств, размер и ассоциативность кэш-памяти всех уровней, эффективность предвыборки данных и команд. Со стороны программного обеспечения IPC определяется характером вычислительной нагрузки: интенсивностью и локальностью обращений к памяти, соотношением арифметических, целочисленных и операций с плавающей точкой, плотностью ветвлений.

Непосредственно изменить IPC ни операционная система, ни пользовательское приложение не могут, поскольку эта величина является производной от физических возможностей ядра и скомпилированного машинного кода. Однако разработчики могут косвенно влиять на него через оптимизацию программ: улучшение локальности данных, уменьшение количества условных переходов и выравнивание кода, что позволяет более полно загрузить вычислительные ресурсы процессора и, как следствие, повысить IPC.

\subsection{Прерывания}

С точки зрения процессора, обработка прерывания является синхронной операцией, включающей сохранение аппаратного контекста, смену режима работы на привилегированный и переход по векторному адресу. Ключевое влияние на загруженность процессора оказывают не сами прерывания, а совокупные затраты времени на выполнение кода их обработчиков.

Для количественной оценки вклада прерываний в общую загрузку процессора может использоваться доля процессорного времени, затрачиваемого на исполнение кода всех обработчиков прерываний:
\begin{equation}
	T_{interrupts} = N_{interrupts} \cdot \overline{t}_{interrupts}
\end{equation}
где:
\begin{enumerate}
	\item $N_{interrupts}$ -- общее количество прерываний;
	\item $\overline{t}_{interrupts}$ -- среднее время выполнения обработчика прерывания.
\end{enumerate}

Тогда относительный вклад в загрузку процессора вычисляется как:
\begin{equation}
	U_{CPU}^{interrupts} = \frac{T_{interrupts}}{T_{total}} \times 100\%
\end{equation}
где:
\begin{enumerate}
	\item $T_{interrupts}$ -- доля времени, затрачиваемого на исполнение кода всех обработчиков прерываний;
	\item $T_{total}$ -- полное время.
\end{enumerate}

\section{Анализ известных способов сбора данных об использовании процессора}

\subsection{proc}

\texttt{proc} -- это виртуальная файловая система, которая используется в качестве интерфейса к структурам ядра операционной системы. Большинство расположенных в ней файлов доступны только для чтения, но некоторые файлы позволяют изменять переменные ядра. 

\newpage

\subsubsection{/proc/stat}

Данный файл содержит статистику системы. Общие записи включают: 
\begin{enumerate}
	\item строка с меткой \texttt{cpu}, содержащая совокупное время, проведенное процессором в различных режимах с момента загрузки системы;
	\item строка с меткой \texttt{intr}, содержащая количество прерываний, случившихся с момента загрузки системы;
	\item строка с меткой \texttt{ctxt}, содержащая количество переключений контекста в системе.
\end{enumerate}

\subsubsection{/proc/cpuinfo}

Данный файл содержит записи информации о процессоре. Общие записи включают:
\begin{enumerate}
	\item строки с меткой \texttt{cpu cores}, содержащие количество физических и логических ядер процессора;
	\item строки с меткой \texttt{cpu MHz}, содержащие текущую тактовую частоту конкретного ядра процессора.
\end{enumerate}

\subsubsection{/proc/interrupts}

Данный файл содержит записи количества аппаратных прерываний с момента запуска системы, соответствующих каждой линии IRQ.

\subsubsection{/proc/softirq}

Данный файл содержит записи количества программных прерываний с момента запуска системы.

\subsection{sysfs}

\texttt{sysfs} -- это виртуальная файловая система, файлы которой содержат информацию об устройствах, модулях ядра, других файловых системах и компонентах ядра.

Директория \texttt{/sys/devices/system/cpu/cpu*/cpufreq} содержит файлы с информацией о тактовой частоте конкретного ядра процессора:
\begin{enumerate}
	\item \texttt{scaling\_cur\_freq} -- файл, содержащий текущее значение тактовой частоты ядра процессора;
	\item \texttt{scaling\_min\_freq} -- файл, содержащий минимальное значение тактовой частоты ядра процессора;
	\item \texttt{scaling\_max\_freq} -- файл, содержащий максимальное значение тактовой частоты ядра процессора;
\end{enumerate}

Из-за отсутствия поддержки управления частотой процессора в ядре или неподходящей аппаратной платформы данная директория может быть недоступна.

Данная файловая система не предоставляет данных об остальных характеристиках в полном объеме, поскольку её архитектура ориентирована на представление иерархии устройств и их параметров, а не агрегированной статистики работы подсистем ядра. 

\subsection{perf\_event.h}

Заголовочный файл \texttt{linux/perf\_event.h} является частью API ядра Linux и предоставляет программный интерфейс для взаимодействия с подсистемой мониторинга производительности. Основное назначение заключается в предоставлении механизмов для создания, управления и чтения событий производительности из пространства ядра. Событие -- это именованный абстрактный ресурс, потребление которого ядром или процессором может быть измерено.

Структура \texttt{struct perf\_event} предназначена для определения собы-тия. Код данной структуры представлен в листинге~\ref{lst:perf_event.c}.

\includelistingpretty{perf_event.c}{c}{Код структуры struct perf\_event}

Структура \texttt{struct perf\_event\_attr} предназначена для определения всех параметров события. Код данной структуры представлен в листинге~\ref{lst:perf_event_attr.c}.

\includelistingpretty{perf_event_attr.c}{c}{Код структуры struct perf\_event\_attr}

Перечислимый тип \texttt{enum perf\_type\_id} представляет тип события. Код данного типа представлен в листинге~\ref{lst:perf_type_id.c}.

\newpage

\includelistingpretty{perf_type_id.c}{c}{Код перечислимого типа enum perf\_type\_id}

Перечислимый тип \texttt{enum perf\_hw\_id} представляет тип аппаратного события. Код данного типа представлен в листинге~\ref{lst:perf_hw_id.c}.

\includelistingpretty{perf_hw_id.c}{c}{Код перечислимого типа enum perf\_hw\_id}

Основным системным вызовом для работы с событиями является \texttt{perf\_event\_open()}. Он  создаёт   файловый   дескриптор,   через  который  можно  получать  измерения производительности. Каждый файловый дескриптор соответствует одному измеряемому  событию. При этом события можно группировать для одновременного измерения. Сигнатура данного системного вызова представлена в листинге~\ref{lst:perf_event_open.c}

\includelistingpretty{perf_event_open.c}{c}{Сигнатура системного вызова perf\_event\_open()}

Интерфейс \texttt{perf\_event.h} не предоставляет исчерпывающих данных о процессорном времени, тактовой частоте и прерываниях, поскольку его архитектурное назначение заключается в мониторинге низкоуровневых событий производительности на основе аппаратных и программных счётчиков, а не в агрегации системной статистики.

\subsection*{Сравнение способов}

Сравнительный анализ способов сбора данных об использовании процессора представлен в таблице~\ref{tbl:methods-compare}.

\begin{table}[H]
	\centering
	\caption{Сравнительный анализ способов сбора данных об использовании процессора}
	\label{tbl:methods-compare}
	\begin{tabularx}{\textwidth}{|X|X|X|X|}
		\hline
		\textbf{Критерий} & \textbf{proc} & \textbf{sysfs} & \textbf{perf\_event.h} \\
		\hline
		Данные о процессорном времени & Да & Нет & Нет \\
		\hline
		Данные о тактовой частоте & Да & Да & Нет \\
		\hline
		Данные о количестве инструкций за такт & Нет & Нет & Да \\
		\hline
		Данные о прерываниях & Да & Нет & Нет \\
		\hline
	\end{tabularx}
\end{table}

\section{Формализация задачи}

Необходимо реализовать загружаемый модуль ядра Linux, предоставляющий информацию из пространства ядра в пространство пользователя об использовании системой за определенный промежуток времени ядер процессора. 